\chapter{Introdução}
\label{cap:cap1}

A região Nordeste do Brasil, caracterizada por sua vulnerabilidade climática, enfrenta historicamente desafios relacionados à escassez hídrica. A gestão eficiente da água não depende apenas do armazenamento, mas também da garantia de integridade das infraestruturas críticas, como barragens e reservatórios. O apoio tecnológico à gestão pode proporcionar maior eficiência na tomada de decisões e mitigação de riscos \cite{hernandez2019}.

No entanto, conforme destacado por \citeonline{khan2024}, a dependência de métodos tradicionais, frequentemente onerosos e de difícil implementação em larga escala, dificulta o monitoramento contínuo dessas estruturas. Essa lacuna tecnológica limita a capacidade de detecção precoce de anomalias, acentuando a vulnerabilidade das barragens a falhas físicas e riscos ambientais iminentes.

O monitoramento estrutural de grandes barragens é um processo complexo e oneroso. Atualmente, a vigilância dessas estruturas no Brasil depende, predominantemente, de inspeções presenciais. Segundo \citeonline{andrade2022}, o método manual de coleta de informações é inadequado por dois fatores principais: primeiro, a dinâmica dos recursos hídricos exige dados atualizados rapidamente para evitar prejuízos; segundo, há um risco significativo à vida humana durante visitas técnicas em locais ermos e de difícil acesso.

O problema central, portanto, reside na dificuldade de realizar um monitoramento frequente, seguro e economicamente viável que englobe tanto variáveis hidrológicas (como nível e vazão) quanto a integridade física da estrutura (como rachaduras e infiltrações). Embora existam soluções comerciais de telemetria, como as plataformas da Trimble (e.g., Trimble 4D Control) e da Bentley Systems (e.g., Sensemetrics), em certos contextos, elas podem esbarrar em limitações técnicas e orçamentárias.

Sistemas de alto desempenho geralmente exigem conectividade robusta para enviar grandes volumes de dados para processamento em nuvem (cloud computing), fazendo com que a maioria das soluções de IoT para barragens foquem na coleta e entrega de dados numéricos, desconsiderando ou pouco abarcando o contexto da análise visual, que depende do processamento de imagens. Barragens como a de Oiticica situam-se em áreas longínquas com conexão de internet limitada (4G instável ou satélite), inviabilizando o envio de streams de vídeo para processamento em tempo real. Além disso, soluções que utilizam apenas sensores físicos pontuais (como extensômetros) podem não cobrir toda a vasta superfície de concreto, deixando fissuras visíveis não instrumentadas sem detecção.

Este artigo propõe o desenvolvimento de um sistema híbrido de monitoramento que une IoT e Inteligência Artificial. A inovação principal consiste na aplicação de Edge Computing (computação na borda) para processamento de imagens. Utilizando hardware embarcado e algoritmos de detecção de objetos, o sistema processa localmente as imagens capturadas, identifica fissuras no concreto e transmite ao servidor remoto apenas metadados (localização e alerta), via protocolo MQTT ou HTTP, contornando a limitação de banda de internet. Adicionalmente, propõe-se a integração desses dados em um dashboard com visualização 3D da barragem, permitindo georreferenciar as falhas estruturais com precisão no modelo virtual.
